\documentclass[11pt,a4paper]{article}
\usepackage[utf8]{inputenc}
\usepackage[T1]{fontenc}
\usepackage{amsmath,amssymb}
\usepackage{booktabs}
\usepackage{multirow}
\usepackage{graphicx}
\usepackage{xcolor}
\usepackage{hyperref}
\usepackage[margin=2.5cm]{geometry}
\usepackage{enumitem}
\usepackage{subcaption}
\usepackage{tablefootnote}

\newcommand{\sys}{OntoCoDec}
\newcommand{\bfval}[1]{\textbf{#1}}

\title{\sys{}: Type-Valid Information Extraction via\\Compiled Ontology Constraints}

\author{Anonymous}
\date{}

\begin{document}
\maketitle

% ============================================================
\begin{abstract}
% ============================================================

Large Language Models (LLMs) produce systematic ontology constraint violations when generating structured Information Extraction (IE) outputs, assigning entity types or argument roles that contradict the schema definition. We present \sys{}, a generalizable framework for ontology-constrained structured generation that automatically compiles IE schemas into two composable decoding constraint layers: a grammar layer enforcing structural validity and type vocabularies via context-free grammars, and a semantic layer enforcing cross-field type constraints (domain/range rules, role restrictions) through stateful logit masking with a dynamic entity registry. The two layers compose via bitwise AND of token-level masks with an enabling dependency: grammar ensures every generated type token belongs to the valid vocabulary, which enables the semantic layer to look up and enforce relational constraints for that type. We evaluate \sys{} on two structurally distinct IE tasks: Relation Extraction (RE) on SciERC and DocRED, and Event Extraction (EE) on WikiEvents, across three model families (Qwen3-8B, Llama-3.1-8B, Mistral-7B). On RE, \sys{} reduces domain/range violations by 44.7--99.4\% with neutral or positive F1 impact. On EE, it eliminates event vocabulary and role vocabulary violations entirely (from 10--89\% to 0\%) across all models, while post-hoc filtering discards 99.6--100\% of arguments. Ablation studies confirm that the grammar layer alone cannot enforce context-dependent constraints (role violations persist at 64--77\%), establishing the necessity of the two-layer architecture.

\end{abstract}

% ============================================================
\section{Introduction}
\label{sec:intro}
% ============================================================

LLM-based information extraction systems increasingly adopt structured generation, prompting models to output JSON or other schema-conforming formats that encode entities, relations, or events. A critical but under-studied failure mode in this paradigm is ontology constraint violation: the model generates structurally well-formed output where the extracted types contradict the ontology definition. For instance, a relation extraction system may assign a \texttt{Method} entity as the tail of a \texttt{Feature-of} relation whose ontology restricts tails to \texttt{OtherScientificTerm}, producing a triple that is syntactically valid JSON but semantically inconsistent with the schema.

These violations are not rare edge cases. Our systematic measurement across two RE benchmarks (SciERC with 7 relation types, DocRED with 96 relation types) reveals per-triple domain/range violation rates of 15.4\% and 18.7\% respectively, with violation rates driven by constraint tightness rather than ontology size. On Event Extraction (WikiEvents, 56 event types, 80 argument roles), role vocabulary violations reach 64--89\% in zero-shot settings, where the combinatorial complexity of per-event-type role constraints far exceeds what models can internalize from prompts alone.

Prompt-level interventions offer partial mitigation. Schema-informed prompting, where explicit constraint descriptions are included in the prompt, reduces violations by approximately 62\% on SciERC but only 32\% on DocRED, confirming that prompt-level approaches scale poorly with ontology complexity. Critically, prompt intervention is unreliable at the per-relation level: for the \texttt{Evaluate-for} relation on SciERC, schema-informed prompting increases the violation rate from 24.1\% to 29.9\%, as the model's distributional priors actively conflict with the stated constraints.

Post-hoc filtering, removing constraint-violating outputs after generation, preserves precision but sacrifices information coverage. On DocRED, post-hoc filtering discards 48.1\% of extracted triples. On WikiEvents EE, the situation is far worse: 99.6--100\% of arguments must be discarded to achieve zero violations, rendering post-hoc filtering non-viable.

We present \sys{}, a framework that automatically compiles IE ontology schemas into decoding-time constraints. The framework consists of two composable layers with an enabling dependency:
(1) a \textbf{grammar layer} that compiles type vocabularies into context-free grammar constraints via XGrammar, ensuring every generated type token belongs to the valid ontology vocabulary;
(2) a \textbf{semantic layer} that maintains a stateful entity registry during generation and applies dynamic logit masking to enforce cross-field constraints (domain/range for RE, role-entity type restrictions for EE).
The grammar layer is a prerequisite for the semantic layer: without guaranteed valid type vocabulary, the semantic layer encounters unknown types and cannot enforce relational constraints. This enabling dependency, confirmed by ablation, distinguishes our architecture from approaches that treat constraint layers as independent.

We classify ontology constraints by prefix-decidability: (a) purely local constraints decidable from vocabulary membership (grammar layer), (b) prefix-decidable constraints requiring generation state (semantic layer with entity registry), and (c) truly non-local constraints like forward coreference that lie outside the scope of forward-only constrained decoding. This taxonomy provides principled guidance for which constraints can be enforced at decoding time and which require alternative mechanisms.

Our contributions are:
\begin{enumerate}[nosep]
\item A systematic measurement of ontology constraint violations in LLM-based IE across two tasks (RE and EE), three model families, and four benchmarks, establishing violation rates, their distribution across constraint types, and the failure modes of prompt-level and post-hoc mitigation.
\item A composable two-layer constrained decoding framework that automatically compiles ontology schemas into grammar and semantic constraints with an enabling dependency, reducing domain/range violations by 44.7--99.4\% on RE and eliminating vocabulary violations on EE, with neutral or positive F1 impact.
\item A decidability taxonomy for IE ontology constraints that delineates which constraint classes are enforceable at decoding time, providing theoretical grounding for the framework's scope and limitations.
\item Cross-task validation demonstrating that the framework generalizes from Relation Extraction to Event Extraction without architectural modification, with ablation confirming the necessity of both layers for context-dependent constraints.
\end{enumerate}

% ============================================================
\section{Background}
\label{sec:background}
% ============================================================

\paragraph{Grammar-Constrained Decoding.}
Grammar-constrained decoding restricts LLM generation to strings accepted by a formal grammar, typically a context-free grammar (CFG). At each decoding step, a token mask is computed from the set of tokens that can extend the current prefix into a valid string. XGrammar provides efficient CFG-based constrained decoding with precomputed token masks. This approach guarantees structural validity (valid JSON, correct field names) and can enforce vocabulary constraints (generated values belong to a predefined set). However, CFGs cannot express constraints that depend on previously generated values within the same output, such as ``the entity type of this relation's head must be compatible with the relation type.''

\paragraph{IE Ontology Structure.}
An IE ontology defines the type system for extracted information. For Relation Extraction, an ontology $\mathcal{O}_{RE} = (\mathcal{E}, \mathcal{R}, \delta)$ consists of entity types $\mathcal{E}$, relation types $\mathcal{R}$, and domain/range constraints $\delta: \mathcal{R} \to 2^{\mathcal{E}} \times 2^{\mathcal{E}}$ specifying valid (head type, tail type) pairs for each relation. For Event Extraction, an ontology $\mathcal{O}_{EE} = (\mathcal{V}, \mathcal{A}, \mathcal{E}, \rho)$ defines event types $\mathcal{V}$, argument roles $\mathcal{A}$, entity types $\mathcal{E}$, and role-entity constraints $\rho: \mathcal{V} \times \mathcal{A} \to 2^{\mathcal{E}}$ specifying which entity types are valid fillers for each role under each event type.

\paragraph{Constraint Tightness.}
We define the tightness of a relation type $r$ as $\tau(r) = 1 - |\delta(r)| / |\mathcal{E}|^2$, measuring how restrictive its domain/range constraint is. A tightness of 1.0 means only one (head, tail) type pair is valid; 0.0 means any pair is allowed. We hypothesize and confirm that violation rates correlate with constraint tightness rather than ontology size.

% ============================================================
\section{Violation Landscape}
\label{sec:landscape}
% ============================================================

Before introducing our constrained decoding framework, we conduct a systematic measurement of ontology constraint violations in LLM-based IE. This landscape study is an independent empirical contribution that motivates the framework design.

\paragraph{Violation Taxonomy.}
We define five violation types across RE and EE:
\begin{itemize}[nosep]
\item \textbf{ET} (Entity Type): generated entity type not in the ontology's valid type set.
\item \textbf{RV} (Relation Vocabulary, RE) / \textbf{EV} (Event Vocabulary, EE): generated relation or event type not in the ontology definition.
\item \textbf{DR} (Domain/Range, RE) / \textbf{RoV} (Role Violation, EE): the type assignment of a relation's arguments or event's argument violates the ontology's type constraints.
\item \textbf{AC} (Argument Constraint, EE): argument entity type not in the allowed set for that (event type, role) pair.
\item \textbf{CF} (Cross-Field): the same entity mention is assigned different types across different relations or events in the same document.
\end{itemize}

\paragraph{RE Violation Rates.}
Table~\ref{tab:re_violations} reports violation rates for unconstrained zero-shot generation on SciERC and DocRED. Key findings: (1) DR rates are comparable across ontology sizes (SciERC 15.4\%, DocRED 18.5\% for Qwen3-8B) despite a 14$\times$ difference in relation count, suggesting violation rates are driven by constraint tightness, not schema complexity. (2) RV rates vary dramatically across models and datasets: Qwen3-8B produces 29.7\% out-of-vocabulary relation names on DocRED but only 1.4\% on SciERC, reflecting the difficulty of adhering to a 96-type vocabulary without explicit enforcement. (3) Schema-informed prompting on SciERC reduces overall DR from 15.4\% to 5.8\% ($-$62\%), but per-relation analysis reveals that three relation types retain $>$10\% violation rates, and \texttt{Evaluate-for} worsens from 24.1\% to 29.9\%.

\begin{table}[t]
\centering
\small
\caption{Ontology constraint violation rates (\%) for unconstrained zero-shot RE. DR = domain/range, RV = relation vocabulary, ET = entity type, CF = cross-field.}
\label{tab:re_violations}
\begin{tabular}{llcccc}
\toprule
Dataset & Model & DR & RV & ET & CF \\
\midrule
\multirow{2}{*}{SciERC} & Qwen3-8B & 15.4 & 1.4 & -- & -- \\
& Llama-3.1-8B & 13.3 & 0.6 & -- & -- \\
\midrule
\multirow{2}{*}{DocRED} & Qwen3-8B & 18.5 & 29.7 & 2.5 & -- \\
& Llama-3.1-8B & 26.1 & 16.0 & -- & -- \\
\bottomrule
\end{tabular}
\end{table}

\paragraph{EE Violation Rates.}
Table~\ref{tab:ee_violations} presents violation rates for zero-shot Event Extraction on WikiEvents (246 documents, KAIROS ontology: 56 event types, 80 argument roles, 27 entity types). Role violations (RoV) dominate at 64--89\%, an order of magnitude higher than RE's domain/range violations. This disparity arises because EE role constraints are conditioned on event type: each of 56 event types permits a different subset of 80 roles, creating a combinatorial constraint space that models cannot internalize from prompts. Event vocabulary violations (EV) range from 10.5\% (Llama) to 37.3\% (Mistral), reflecting varying capacity to reproduce hierarchical event type names. Entity F1 is 0\% across all models in unconstrained mode because the generated JSON structure does not conform to the expected parser format, preventing entity mention identification.

\begin{table}[t]
\centering
\small
\caption{Ontology constraint violation rates (\%) for unconstrained zero-shot EE on WikiEvents. EV = event vocabulary, RoV = role vocabulary, AC = argument constraint.}
\label{tab:ee_violations}
\begin{tabular}{lccccc}
\toprule
Model & Parse\% & EV & RoV & AC & CF \\
\midrule
Qwen3-8B & 88.6 & 24.0 & 74.9 & 0.3 & 0.0 \\
Llama-3.1-8B & 91.9 & 10.5 & 88.8 & 0.1 & 0.0 \\
Mistral-7B & 78.9 & 37.3 & 63.9 & 0.0 & 0.0 \\
\bottomrule
\end{tabular}
\end{table}

% ============================================================
\section{\sys{} Framework}
\label{sec:method}
% ============================================================

\sys{} takes an IE ontology schema as input and automatically compiles it into two composable constraint layers applied during autoregressive decoding.

\subsection{Grammar Layer}
\label{sec:grammar}

The grammar layer compiles the ontology's type vocabulary into a CFG that constrains generation to produce valid JSON with type tokens drawn exclusively from the ontology definition. For RE, this means relation names $\in \mathcal{R}$ and entity type names $\in \mathcal{E}$. For EE, event type names $\in \mathcal{V}$, argument role names $\in \mathcal{A}$, and entity type names $\in \mathcal{E}$. The CFG is compiled via XGrammar into precomputed token masks $M_{\text{grammar}}(t)$ at each position $t$.

For EE, hierarchical event type names (e.g., \texttt{Conflict.Attack.DetonateExplode}) create first-token collisions in BPE tokenization. We resolve this by flattening to leaf-node names, reducing collision rate from 96.1\% to 3.9\%.

\subsection{Semantic Layer}
\label{sec:semantic}

The semantic layer maintains a dynamic entity registry $\mathcal{R}_t$ that tracks entity mentions and their assigned types as they are generated. When the model generates a type-constrained token (e.g., the entity type for a relation's head argument), the semantic layer computes a mask $M_{\text{semantic}}(t)$ based on the current context: (1) for RE, if the relation type $r$ has been generated, only entity types in $\delta(r)_{\text{head}}$ (or $\delta(r)_{\text{tail}}$) are allowed; (2) for EE, if the event type $v$ and role $a$ are known, only entity types in $\rho(v, a)$ are allowed.

The entity registry enables \textbf{continuation masking}: when a multi-token entity mention has been partially generated and matched to a known entity, the semantic layer constrains subsequent tokens to complete that entity's previously assigned type. This is essential for EE where the same entity may appear across multiple events.

\subsection{Composition and Graceful Degradation}
\label{sec:composition}

The two layers compose via bitwise AND: $M(t) = M_{\text{grammar}}(t) \wedge M_{\text{semantic}}(t)$. This composition has an enabling dependency: the grammar layer guarantees that every generated type token is in the valid vocabulary, which enables the semantic layer to look up constraint rules for that type. Without the grammar layer, the semantic layer encounters out-of-vocabulary types and cannot enforce constraints, explaining why a semantic-only configuration achieves only 24\% DR reduction compared to the full system's 94\%.

When $M_{\text{grammar}}(t) \wedge M_{\text{semantic}}(t) = \emptyset$ (the intersection is empty, meaning no token satisfies both constraint layers simultaneously), the system falls back to grammar-only constraints: $M(t) = M_{\text{grammar}}(t)$. This ensures structural validity is maintained as a guaranteed floor even when semantic constraints are too tight for the current generation context. The fallback frequency correlates with constraint tightness: on DocRED ($\bar{\tau}=0.842$, 96 relations), approximately 57 fallbacks per document occur compared to approximately 11 per document on SciERC (7 relations, lower average tightness). This fallback rate directly explains residual DR after constrained decoding: positions where semantic constraints could not be enforced revert to grammar-only control, which has a known DR floor.

\subsection{Decidability Taxonomy}
\label{sec:decidability}

We classify ontology constraints by their prefix-decidability during autoregressive generation:

\begin{enumerate}[nosep,label=(\alph*)]
\item \textbf{Purely local}: decidable from vocabulary membership alone (e.g., entity type $\in \mathcal{E}$). Enforced by the grammar layer.
\item \textbf{Prefix-decidable with state}: decidable given previously generated content and a state tracker (e.g., relation argument types compatible with the generated relation type). Enforced by the semantic layer with the entity registry.
\item \textbf{Truly non-local}: require information not yet generated (e.g., forward coreference, where a later mention determines the type of an earlier one). Outside the scope of single-pass forward decoding.
\end{enumerate}

Cross-field consistency (CF) violations fall into category (c): ensuring the same entity receives the same type across independently generated relations or events requires global coordination that forward-only decoding cannot provide. We discuss CF as an explicit limitation in Section~\ref{sec:analysis}.

% ============================================================
\section{Experiments}
\label{sec:experiments}
% ============================================================

\subsection{Setup}
\label{sec:setup}

\paragraph{Datasets.}
For RE: SciERC (551 abstracts, 7 relation types, 6 entity types) and DocRED (998 documents, 96 relation types, 7 entity types). For EE: WikiEvents (246 documents, KAIROS ontology with 56 event types, 80 argument roles, 27 entity types, and 245 explicit (event\_type, role) $\to$ allowed\_entity\_types constraints).

\paragraph{Models.}
Qwen3-8B, Llama-3.1-8B-Instruct, Mistral-7B-Instruct-v0.3 for main experiments. Llama-3.1-70B-AWQ-INT4 for scale analysis.

\paragraph{Configurations.}
We evaluate: (1) unconstrained (baseline), (2) grammar-only (grammar layer without semantic), (3) full constrained (\sys{}, both layers), and (4) post-hoc filtering (unconstrained generation followed by removal of violating outputs). All experiments use zero-shot prompting.

\paragraph{Metrics.}
Violation rates (DR, RV, ET, EV, RoV, AC, CF) measured per-triple or per-argument. Task performance via F1 (relation F1 for RE, trigger/entity/argument F1 for EE). Statistical significance via 10,000-iteration bootstrap resampling.

\subsection{Relation Extraction Results}
\label{sec:re_results}

Table~\ref{tab:re_main} presents the main RE results. \sys{} achieves substantial DR reduction across both datasets and models:

\begin{itemize}[nosep]
\item \textbf{SciERC}: Qwen DR 15.4\%$\to$0.93\% ($\downarrow$94\%), Llama DR 13.3\%$\to$0.08\% ($\downarrow$99.4\%). Relation F1 is neutral: Qwen +0.13pp ($p=0.58$), Llama $-$0.14pp ($p=0.011$, negligible magnitude).
\item \textbf{DocRED}: Qwen DR 18.5\%$\to$10.2\% ($\downarrow$44.7\%), Llama DR 26.1\%$\to$2.55\% ($\downarrow$90.2\%). Relation F1 significantly improved: Qwen +0.50pp ($p<0.001$, 95\% CI [+0.23, +0.81]), Llama +0.29pp ($p=0.013$). Entity F1: Qwen +1.23pp ($p<0.001$), Llama $-$0.64pp ($p=0.022$).
\end{itemize}

The cross-model divergence on DocRED reveals an important interaction: Qwen's higher RV baseline (29.7\% vs Llama's 16.0\%) means the grammar layer forces more out-of-vocabulary tokens onto valid relation names, exposing latent DR violations that were previously masked by vocabulary errors. The resulting DR reduction ceiling is determined jointly by constraint tightness and the model's vocabulary compliance baseline, not by tightness alone.

Grammar-only ablation on DocRED confirms the enabling dependency: grammar eliminates RV (29.7\%$\to$0.03\%) and ET (2.5\%$\to$0\%), but DR increases to 26.2\% (above the unconstrained 18.5\%). This latent violation surfacing occurs because grammar forces valid relation names, exposing DR violations that were hidden when the relation name itself was invalid. The semantic layer then reduces DR from 26.2\% to 10.2\%.

\begin{table}[t]
\centering
\small
\caption{RE main results: violation rates (\%) and F1 (\%). Best violation rate per column in \textbf{bold}. $\Delta$F1 shows change from unconstrained.}
\label{tab:re_main}
\begin{tabular}{llcccccc}
\toprule
Dataset & Config & DR & RV & ET & CF & R-F1 & E-F1 \\
\midrule
\multirow{3}{*}{\rotatebox{0}{SciERC}} & Qwen Uncon. & 15.4 & 1.4 & -- & -- & 8.23 & 43.7 \\
& Qwen \sys{} & \bfval{0.93} & \bfval{0.0} & -- & 9.9 & 8.36 & 43.7 \\
& Llama Uncon. & 13.3 & 0.6 & -- & -- & 8.50 & -- \\
& Llama \sys{} & \bfval{0.08} & \bfval{0.0} & -- & 6.6 & 8.36 & -- \\
\midrule
\multirow{4}{*}{\rotatebox{0}{DocRED}} & Qwen Uncon. & 18.5 & 29.7 & 2.5 & -- & 7.74 & 58.1 \\
& Qwen \sys{} & 10.2 & \bfval{0.03} & \bfval{0.0} & 2.6 & 8.24 & 59.4 \\
& Llama Uncon. & 26.1 & 16.0 & -- & -- & -- & -- \\
& Llama \sys{} & \bfval{2.55} & \bfval{0.0} & -- & 12.0 & -- & -- \\
\bottomrule
\end{tabular}
\end{table}

\paragraph{Post-hoc Filtering Comparison.}
Post-hoc filtering removes all triples containing any violation. On SciERC, this discards 16.8\% of triples; on DocRED, 48.1\%. Post-hoc achieves higher precision (by removing incorrect triples) but lower recall: on DocRED, 127 correctly extracted triples (11.2\% of true positives) are discarded because they contain entity type violations while being otherwise correct. \sys{} prevents these violations at generation time, preserving these true positives. Post-hoc also cannot resolve entity type violations (ET residual 0.4\% on SciERC, 2.5\% on DocRED), as it only removes triples without correcting entity lists.

\subsection{Event Extraction Results}
\label{sec:ee_results}

Table~\ref{tab:ee_main} presents the main EE results on WikiEvents. \sys{} eliminates event vocabulary violations (EV) and role vocabulary violations (RoV) entirely across all three models, from 10--37\% and 64--89\% respectively to 0\%. This is the strongest demonstration of the framework's constraint enforcement capability: the combinatorial space of 56 event types $\times$ 80 roles makes these constraints impossible to enforce through prompting alone.

Entity F1 recovers from 0\% to 11.7--17.6\%. This recovery is primarily a format compliance effect: grammar enforcement produces schema-valid JSON that the evaluation parser can process, enabling entity mention identification that was impossible with unconstrained output. Trigger F1 improves by 1.5--1.8pp across models ($p<0.05$ for all; Qwen $p=0.009$, Llama $p=0.018$, Mistral $p=0.045$) due to elimination of out-of-vocabulary event types that were counted as false positives.

\begin{table}[t]
\centering
\small
\caption{EE results on WikiEvents (246 docs): violation rates (\%) and F1 (\%). Uncon.\ = unconstrained, Full = \sys{} (grammar + semantic).}
\label{tab:ee_main}
\begin{tabular}{llccccccc}
\toprule
Model & Config & EV & RoV & AC & CF & Trig-F1 & Ent-F1 \\
\midrule
\multirow{2}{*}{Qwen} & Uncon. & 24.0 & 74.9 & 0.3 & 0.0 & 11.8 & 0.0 \\
& Full & \bfval{0.0} & \bfval{0.0} & 2.1 & 21.1 & 13.4 & 14.1 \\
\midrule
\multirow{2}{*}{Llama} & Uncon. & 10.5 & 88.8 & 0.1 & 0.0 & 5.7 & 0.0 \\
& Full & \bfval{0.0} & \bfval{0.0} & 3.8 & 17.2 & 7.5 & 17.6 \\
\midrule
\multirow{2}{*}{Mistral} & Uncon. & 37.3 & 63.9 & 0.0 & 0.0 & 2.4 & 0.0 \\
& Full & \bfval{0.0} & \bfval{0.0} & 1.3 & 14.9 & 3.9 & 11.7 \\
\bottomrule
\end{tabular}
\end{table}

\paragraph{Post-hoc Filtering on EE.}
Post-hoc filtering is non-viable for EE. Table~\ref{tab:posthoc_ee} shows that 99.6--100\% of arguments must be discarded to eliminate violations, because RoV affects 64--89\% of arguments and EV affects 10--37\% of event instances, creating cascading rejection. This contrasts sharply with RE, where post-hoc discards 17--48\% of triples, and establishes that constrained decoding is the only feasible approach for ontology compliance in complex IE tasks.

\begin{table}[t]
\centering
\small
\caption{Post-hoc filtering comparison: RE vs.\ EE. On EE, near-total argument rejection makes post-hoc non-viable.}
\label{tab:posthoc_ee}
\begin{tabular}{llccc}
\toprule
Task & Model & Args/Triples & Surviving & Rejection\% \\
\midrule
RE (SciERC) & Qwen & -- & -- & 16.8 \\
RE (DocRED) & Qwen & -- & -- & 48.1 \\
\midrule
\multirow{3}{*}{EE (WikiEvents)} & Qwen & 5{,}467 & 23 & 99.6 \\
& Llama & 5{,}053 & 3 & 99.9 \\
& Mistral & 6{,}034 & 0 & 100.0 \\
\bottomrule
\end{tabular}
\end{table}

\subsection{Ablation: Grammar-Only on EE}
\label{sec:ablation_ee}

Table~\ref{tab:ee_ablation} presents the grammar-only ablation for all three models on WikiEvents. The grammar layer eliminates EV entirely (EV$\to$0\% for all models) but RoV remains at 64--77\%, nearly unchanged from unconstrained baselines. This confirms that role vocabulary constraints, which are conditioned on event type context, cannot be expressed in a CFG and require the stateful semantic layer.

AC rises under grammar-only (from 0--0.3\% to 8--14\%), exhibiting the same latent violation surfacing pattern observed in RE: by forcing valid event types, the grammar layer enables AC checking on arguments that were previously attached to invalid events, exposing pre-existing constraint violations. The semantic layer then resolves these violations, reducing AC to 1.3--3.8\%.

Entity F1 under grammar-only (11.9--16.9\%) is comparable to the full system (11.7--17.6\%), confirming that format compliance recovery is driven by the grammar layer's JSON structure enforcement, with the semantic layer contributing constraint satisfaction rather than extraction quality.

\begin{table}[t]
\centering
\small
\caption{EE grammar-only ablation on WikiEvents. Grammar eliminates EV but RoV persists at 64--77\%, proving the semantic layer is necessary for context-dependent constraints.}
\label{tab:ee_ablation}
\begin{tabular}{llccccc}
\toprule
Model & Config & EV & RoV & AC & CF & Ent-F1 \\
\midrule
\multirow{2}{*}{Qwen} & Grammar & 0.0 & 71.5 & 8.2 & 0.6 & 13.5 \\
 & Full & 0.0 & \bfval{0.0} & 2.1 & 21.1 & 14.1 \\
\midrule
\multirow{2}{*}{Llama} & Grammar & 0.0 & 76.8 & 11.5 & 0.4 & 16.9 \\
 & Full & 0.0 & \bfval{0.0} & 3.8 & 17.2 & 17.6 \\
\midrule
\multirow{2}{*}{Mistral} & Grammar & 0.0 & 64.0 & 13.8 & 1.1 & 11.9 \\
 & Full & 0.0 & \bfval{0.0} & 1.3 & 14.9 & 11.7 \\
\bottomrule
\end{tabular}
\end{table}

\subsection{Quantitative Analysis}
\label{sec:analysis}

\paragraph{Cross-Field Violations.}
CF violations are introduced by constrained decoding: unconstrained CF is near-zero (0--0.1\%) while \sys{} CF reaches 9.9\% on SciERC RE and 14.9--21.1\% on WikiEvents EE. CF is entirely attributable to the semantic layer (grammar-only CF $<$1.1\% across all configurations). The mechanism is straightforward: when the semantic layer enforces per-relation or per-event type constraints independently, the same entity can receive different types in different relational contexts.

Root cause analysis on 15 representative EE cases reveals that CF is dominated by ontology-inherent type boundary ambiguity. The top confusion patterns are \texttt{SideOfConflict}$\leftrightarrow$\texttt{Sensor} (3 cases, surface-form overlap in military context), \texttt{Person}$\leftrightarrow$\texttt{CommercialOrganization} (2 cases, named entities ambiguous between individual and corporate reference), and \texttt{GPE}$\leftrightarrow$\texttt{Location} (1 case). The dominant CF pattern involves the same entity (e.g., ``the US'') being assigned \texttt{GPE} in one event and \texttt{Government} in another, where the per-event context legitimately determines different type assignments.

CF violations produce individually type-valid triples (each triple satisfies its own domain/range or role constraints), unlike DR/RoV violations which produce structurally invalid output. We consider CF a soft consistency preference rather than a hard constraint, and classify it under the ``truly non-local'' category (c) of our decidability taxonomy. Resolving CF requires a global entity-type consensus mechanism that operates across independently generated relations/events, which we identify as future work.

\paragraph{Enabling Dependency Quantification.}
The enabling dependency between grammar and semantic layers is quantified by comparing grammar-only and full system performance. On WikiEvents, grammar alone achieves EV$\to$0\% but leaves RoV at 64--77\%. Adding the semantic layer reduces RoV to 0\%. This clean separation, where grammar handles vocabulary (class a constraints) and semantic handles context-dependent rules (class b constraints), validates the decidability taxonomy.

On DocRED, the latent violation surfacing effect provides additional evidence: grammar-only DR (26.2\%) exceeds unconstrained DR (18.5\%) because grammar forces valid relation names, enabling DR measurement on triples that previously had invalid relation types. The semantic layer then reduces this elevated DR to 10.2\%.

\paragraph{Residual AC in EE.}
Residual AC violations (1.3--3.8\%) after full constrained decoding reflect the model's semantic confusion between ontologically similar entity types. BPE collision analysis confirms that Qwen (152K vocab) and Llama (128K vocab) have zero first-token collisions among the 27 KAIROS entity types, so their AC residual (2.1\%, 3.8\%) is purely semantic. Mistral (SentencePiece, 32K vocab) has 2 collision groups (\texttt{SID/SEN/SIDE}$\to$\texttt{S}, \texttt{BAL/BOD}$\to$\texttt{B}); its lower AC residual (1.3\%) reflects partial collision mitigation.

\subsection{Supplementary Experiments}
\label{sec:supplementary}

\paragraph{Violations Persist at Scale but Diminish.}
Table~\ref{tab:70b} presents results for Llama-3.1-70B-AWQ-INT4 on SciERC and DocRED. At 70B parameters, unconstrained violations decrease substantially compared to 8B models (SciERC DR 6.5\% vs 13--15\%, DocRED DR 5.8\% vs 18--26\%), confirming that larger models exhibit stronger schema adherence. However, violations do not vanish: 6.5\% DR on SciERC means approximately 1 in 15 triples is structurally invalid, and 5.8\% on DocRED means 1 in 17.

\sys{} at 70B achieves DR$\to$0\% on SciERC (complete elimination) and DR$\to$1.1\% on DocRED (81\% reduction). R-F1 impact is negligible: +0.02pp on SciERC, +0.2pp on DocRED. These results indicate that while scale reduces the violation rate, constrained decoding provides the formal guarantee that scale alone cannot. For compliance-critical applications such as knowledge graph construction in regulated domains, even reduced violation rates may be unacceptable without enforcement.

\begin{table}[t]
\centering
\small
\caption{Scale analysis: Llama-3.1-70B on SciERC and DocRED. Violations diminish with scale but do not vanish; \sys{} provides formal guarantees.}
\label{tab:70b}
\begin{tabular}{llcccccc}
\toprule
Dataset & Config & DR & RV & ET & CF & R-F1 & E-F1 \\
\midrule
\multirow{2}{*}{SciERC} & 70B Uncon. & 6.5 & 0.6 & 0.0 & 0.0 & 11.22 & 35.17 \\
& 70B \sys{} & \bfval{0.0} & \bfval{0.0} & 0.0 & 3.7 & 11.24 & 35.14 \\
\midrule
\multirow{2}{*}{DocRED} & 70B Uncon. & 5.8 & 3.4 & 0.05 & 0.08 & 10.6 & 63.2 \\
& 70B \sys{} & \bfval{1.1} & \bfval{0.0} & \bfval{0.0} & 4.5 & 10.8 & 64.0 \\
\bottomrule
\end{tabular}
\end{table}

\paragraph{Fine-Tuned Model Composability.}
To verify that \sys{} composes with task-specific fine-tuning, we evaluate a fine-tuned Qwen3-8B on WikiEvents EE. Fine-tuning dramatically reduces baseline violations (EV 24.0\%$\to$0.3\%, RoV 74.9\%$\to$2.8\%) and improves task performance (Trigger F1 11.8\%$\to$41.0\%, Argument F1 0.09\%$\to$27.5\%, Entity F1 0\%$\to$54.0\%). Applying \sys{} to the fine-tuned model eliminates residual EV and RoV entirely, further reduces AC from 1.5\% to 0.3\%, and yields modest F1 gains (Trigger +0.5pp, Argument +0.3pp, Entity +1.2pp). The 95\% DR reduction from the already-low fine-tuned baseline demonstrates that \sys{} provides complementary value even when models have been trained on task-specific data. CF rises from 0.1\% to 1.8\% under constrained decoding, consistent with the pattern observed in zero-shot settings but at a much lower magnitude.

\paragraph{Latency Overhead.}
\sys{} introduces 13.1\% mean latency overhead on SciERC (4.16s$\to$4.71s per document) and 21.5\% on DocRED (8.8s$\to$10.7s). The overhead scales with ontology complexity: more relation types increase semantic mask lookup time, and longer documents increase entity registry maintenance cost. P95 per-document overhead ratios are 1.26 (SciERC) and 1.42 (DocRED). WikiEvents EE overhead is estimated at approximately 26\%, reflecting the largest constraint space (56$\times$80$\times$27). These overheads are acceptable for batch processing and can be reduced through constraint caching and batch constraint evaluation.

\paragraph{Bootstrap Confidence Intervals for EE.}
Entity F1 improvements are significant across all models ($p<0.001$; Qwen 95\% CI [+11.0, +17.2]pp, Llama [+14.2, +21.1]pp, Mistral [+8.6, +14.8]pp), though this significance reflects format compliance recovery rather than extraction quality improvement. Trigger F1 improvements are significant for Qwen ($p=0.009$, CI [+0.4, +2.7]pp) and Llama ($p=0.018$, CI [+0.3, +3.2]pp), with Mistral borderline ($p=0.045$, CI [+0.0, +2.9]pp). The consistent positive direction across three models supports the claim that vocabulary violation elimination reduces false positive event detections.

% ============================================================
\section{Related Work}
\label{sec:related}
% ============================================================

\paragraph{Constrained Decoding for Structured Generation.}
Grammar-constrained decoding enforces structural validity through formal grammars. Guidance, Outlines, and XGrammar compile regular expressions or CFGs into token masks for JSON schema compliance. LMQL extends this with scripted constraints. These approaches guarantee syntactic validity but cannot express semantic constraints that depend on previously generated values. \sys{} builds on grammar-constrained decoding (using XGrammar) and adds a semantic layer for cross-field type constraints.

\paragraph{Type-Constrained Generation.}
Type-checking approaches for code generation, such as those compiling type systems into decoding constraints, enforce that generated programs are well-typed. \sys{} applies an analogous principle to IE ontologies, where ``well-typed'' means entity and relation types satisfy domain/range constraints. The key difference is that IE ontologies define a richer constraint structure with cross-field dependencies (an entity's valid types depend on which relation it participates in), requiring the stateful entity registry that distinguishes our semantic layer.

\paragraph{LLM-Based Information Extraction.}
Recent work frames IE as structured generation: models produce JSON-formatted entity, relation, or event structures. GoLLIE, CodeIE, and InstructUIE demonstrate strong zero-shot and few-shot extraction using schema-guided prompting. These approaches assume the model will respect ontology constraints given sufficient prompt engineering. Our violation landscape study (Section~\ref{sec:landscape}) demonstrates that this assumption fails systematically, motivating decoding-time enforcement.

\paragraph{Event Extraction Methods.}
Event extraction has evolved from pipeline approaches (separate trigger detection and argument extraction) to end-to-end generation. DEGREE, CEDAR, and Text2Event cast EE as conditional generation with schema-aware prompting. WikiEvents and the KAIROS ontology provide a benchmark with rich type constraints (56 event types, 80 roles, 27 entity types). Our work is complementary: \sys{} can be applied on top of any generation-based EE system to enforce ontology compliance.

\paragraph{Post-hoc Validation.}
An alternative to constrained decoding is generating freely and filtering or correcting violations afterward. Self-consistency and self-refinement approaches re-prompt the model to fix errors. Our post-hoc filtering experiments (Section~\ref{sec:ee_results}) demonstrate that this strategy fails catastrophically for complex IE tasks where violation rates are high, discarding up to 100\% of extracted arguments.

% ============================================================
\section{Conclusion}
\label{sec:conclusion}
% ============================================================

We presented \sys{}, a generalizable framework for ontology-constrained information extraction that automatically compiles IE schemas into composable grammar and semantic decoding constraints. Through systematic evaluation on two structurally distinct IE tasks (Relation Extraction and Event Extraction) across three model families and four benchmarks, we demonstrated that: (1) LLM-based IE systems produce substantial ontology violations (15--89\% depending on task and constraint type), (2) prompt-level and post-hoc mitigation are insufficient or non-viable for complex ontologies, (3) the two-layer architecture with enabling dependency provides principled constraint enforcement with neutral or positive F1 impact, and (4) the framework generalizes across tasks without architectural modification.

\paragraph{Limitations.}
Cross-field consistency violations (CF) are introduced by per-relation/per-event independent type enforcement and range from 2.6\% to 21.1\% across configurations. Our decidability taxonomy classifies CF as truly non-local (class c), outside the scope of forward-only constrained decoding. Resolving CF requires a global entity-type consensus mechanism, which we leave to future work.

Our evaluation is limited to zero-shot and single fine-tuned settings. While the framework is model-agnostic (validated on three 8B-class models and one 70B model with zero code modification), we have not evaluated on proprietary models where decoding-time intervention may not be available.

The violation landscape study and the framework are validated on two IE tasks (RE and EE). Generalization to other structured prediction tasks with ontology constraints (e.g., NER with type hierarchies, slot filling, semantic role labeling) is a natural extension that we identify as future work.

\paragraph{Future Work.}
Three directions emerge: (1) a global entity-type consensus pass before per-relation enforcement to resolve CF violations, (2) extension to additional IE tasks (NER, slot filling, semantic role labeling) to further validate framework generality, and (3) integration with beam search decoding, where preliminary results suggest $k=4$ can simultaneously reduce DR and CF through diverse candidate selection.

\end{document}
